\section{习题课 \#3（9月27日）}

\subsection{线性相关、线性无关与交换定理}

\begin{remark}
当 $n\geq2$ 时，$S_n$ 中奇排列与偶排列数目相同。一个简洁证明是
\begin{equation}
 \sum_{\sigma\in S_n}\sgn(\sigma)
 =\det\begin{pmatrix}
 1&\cdots&1\\
 \vdots&&\vdots\\
 1&\cdots&1
 \end{pmatrix}=0.
\end{equation}
\end{remark}

\begin{theorem}[方阵列向量线性无关的等价条件]
设 $A=(\alpha_1,\ldots,\alpha_n)\in\F^{n\times n}$。下列条件等价：
\begin{enumerate}
  \item $\alpha_1,\ldots,\alpha_n$ 线性无关；
  \item $\sum_{j=1}^n c_j\alpha_j=0$ 只有零解；
  \item $\det A\neq0$；
  \item 对任意 $b\in\F^n$，方程 $Ax=b$ 有唯一解；
  \item $A$ 可逆。
\end{enumerate}
\end{theorem}

\begin{remark}
对于 $s\leq n$ 个向量 $\alpha_1,\ldots,\alpha_s\in\F^n$，线性无关仍等价于齐次关系只有零解，但不再能直接使用一个 $s\times s$ 行列式，除非选取适当的 $s$ 行形成非零子式。
\end{remark}

\begin{proposition}[用线性组合替换一个向量]
设 $\alpha_1,\ldots,\alpha_s$ 线性无关，并令
\begin{equation}
 \beta=\sum_{i=1}^s k_i\alpha_i.
\end{equation}
若 $k_j\neq0$，则
\begin{equation}
 \alpha_1,\ldots,\alpha_{j-1},\beta,
 \alpha_{j+1},\ldots,\alpha_s
\end{equation}
仍线性无关。
\end{proposition}

\begin{proof}
若
\begin{equation}
 r\beta+\sum_{i\neq j}r_i\alpha_i=0,
\end{equation}
代入 $\beta$ 后，$\alpha_j$ 的系数为 $rk_j$。由原向量组线性无关且 $k_j\neq0$，先得 $r=0$，再得所有 $r_i=0$。
\end{proof}

\begin{theorem}[Steinitz 交换定理]
设 $\alpha_1,\ldots,\alpha_s$ 张成向量空间 $V$，而 $\beta_1,\ldots,\beta_t$ 线性无关。则 $t\leq s$，并且可以从 $\alpha_1,\ldots,\alpha_s$ 中删去 $t$ 个向量，使
\begin{equation}
 \beta_1,\ldots,\beta_t
\end{equation}
与余下的 $s-t$ 个 $\alpha$ 向量仍张成 $V$。若原来的 $\alpha$ 向量组是一组基，则替换后仍是一组基。
\end{theorem}

\begin{proof}
对 $t$ 归纳。因 $\beta_1\in\Span(\alpha_1,\ldots,\alpha_s)$，其表示中至少有一个系数非零。由上一命题可用 $\beta_1$ 替换相应的 $\alpha$ 而不改变张成空间。随后将 $\beta_2$ 写成新生成组的线性组合。若 $\beta_2$ 在 $\beta_1$ 之外的所有系数均为零，则 $\beta_1,\beta_2$ 相关，矛盾；故还可替换另一个 $\alpha$。如此继续即可。
\end{proof}

\begin{remark}
有限向量组可按其张成空间比较：
\begin{equation}
 A\preccurlyeq B
 \quad\Longleftrightarrow\quad
 \Span A\subseteq\Span B.
\end{equation}
这是一个预序。将张成同一子空间的向量组视为等价后，所得偏序正是有限生成子空间按包含关系构成的偏序。
\end{remark}

\subsection{秩与线性关系的支撑}

\begin{proposition}
设 $\alpha_1,\ldots,\alpha_s\in\F^n$，并固定 $r\leq\min\{n,s\}$。下列条件等价：
\begin{enumerate}
  \item 任取其中 $r$ 个向量均线性无关；
  \item 每个非平凡线性关系
  \begin{equation}
   \sum_{j=1}^s x_j\alpha_j=0
  \end{equation}
  至少含有 $r+1$ 个非零系数。
\end{enumerate}
\end{proposition}

\begin{proof}
若存在只含不超过 $r$ 个非零系数的非平凡关系，则把相应向量补足到 $r$ 个便得到相关组。反之，若某 $r$ 个向量相关，则其非平凡关系扩充为整个向量组上的关系，非零系数不超过 $r$。
\end{proof}

\begin{proposition}[严格对角占优判别]
设 $A=(a_{ij})\in\F^{n\times n}$，并假设绝对值来自 $\R$ 或 $\C$。若
\begin{equation}
 |a_{ii}|>\sum_{j\neq i}|a_{ij}|,
 \qquad i=1,\ldots,n,
\end{equation}
则 $A$ 可逆。
\end{proposition}

\begin{proof}
若 $Ax=0$ 且 $x\neq0$，取 $k$ 使 $|x_k|=\max_j|x_j|>0$。第 $k$ 行给出
\begin{equation}
 |a_{kk}|\,|x_k|
 \leq\sum_{j\neq k}|a_{kj}|\,|x_j|
 \leq\left(\sum_{j\neq k}|a_{kj}|\right)|x_k|,
\end{equation}
与严格不等式矛盾。
\end{proof}

\subsection{若干行列式恒等式}

\begin{exercise}[行和与列差]
设 $A=(a_{ij})\in\F^{n\times n}$，并令
\begin{equation}
 s_i=\sum_{k=1}^n a_{ik},
 \qquad
 B=(s_i-a_{ij})_{i,j=1}^n.
\end{equation}
证明
\begin{equation}
 \det B=(-1)^n(1-n)\det A.
\end{equation}
\end{exercise}

\begin{proof}
记 $J=\ones\ones^T$。因为 $A\ones=(s_1,\ldots,s_n)^T$，故
\begin{equation}
 B=A(J-I_n).
\end{equation}
矩阵 $J-I_n$ 在 $\Span\{\ones\}$ 上的特征值为 $n-1$，在其余 $n-1$ 维子空间上的特征值为 $-1$。于是
\begin{equation}
 \det(J-I_n)=(n-1)(-1)^{n-1}=(-1)^n(1-n),
\end{equation}
结论随即成立。
\end{proof}

\begin{exercise}[列差行列式与余子式之和]
设 $A=(a_{ij})\in\F^{n\times n}$，列向量为 $a_1,\ldots,a_n$。证明
\begin{equation}
 \det(a_1-a_2,a_2-a_3,\ldots,a_{n-1}-a_n,\ones)
 =\sum_{i,j=1}^n C_{ij}(A).
\end{equation}
\end{exercise}

\begin{proof}
由矩阵行列式引理的多项式形式，
\begin{equation}
 \det(A+t\ones\ones^T)
 =\det A+t\ones^T\adj(A)\ones
 =\det A+t\sum_{i,j}C_{ij}(A).
\end{equation}
另一方面，对列作连续差分可把 $A+t\ones\ones^T$ 的一次项系数化为题设左端行列式，故二者相等。也可直接反复使用列的线性性展开。
\end{proof}

\begin{exercise}[正弦和矩阵的秩]
计算
\begin{equation}
 D_n=\det\bigl(\sin(\alpha_k+\beta_\ell)\bigr)_{k,\ell=1}^n.
\end{equation}
\end{exercise}

\begin{solution}
由加法公式
\begin{equation}
 \sin(\alpha_k+\beta_\ell)
 =\sin\alpha_k\cos\beta_\ell+
   \cos\alpha_k\sin\beta_\ell,
\end{equation}
整个矩阵是两个秩一矩阵之和，故秩至多为 $2$。因此
\begin{equation}
 D_n=0,
 \qquad n\geq3.
\end{equation}
低阶情形为
\begin{equation}
 D_1=\sin(\alpha_1+\beta_1),
\end{equation}
以及
\begin{equation}
 D_2
 =\sin(\alpha_1-\alpha_2)
  \sin(\beta_2-\beta_1).
\end{equation}
原稿也用复指数展开说明：每一列均落在由
$(\e^{\ii\alpha_k})_k$ 与 $(\e^{-\ii\alpha_k})_k$ 张成的二维空间中。
\end{solution}

\begin{exercise}[$\mu$-循环行列式]
设
\begin{equation}
 A=\begin{pmatrix}
 c_0&c_1&\cdots&c_{n-1}\\
 \mu c_{n-1}&c_0&\cdots&c_{n-2}\\
 \vdots&\vdots&\ddots&\vdots\\
 \mu c_1&\mu c_2&\cdots&c_0
 \end{pmatrix},
 \qquad
 f(z)=\sum_{j=0}^{n-1}c_jz^j.
\end{equation}
计算 $\det A$。
\end{exercise}

\begin{solution}
若 $\xi^n=\mu$，令
\begin{equation}
 v(\xi)=(1,\xi,\ldots,\xi^{n-1})^T.
\end{equation}
直接计算得
\begin{equation}
 Av(\xi)=f(\xi)v(\xi).
\end{equation}
当 $\mu\neq0$ 时，方程 $\xi^n=\mu$ 的 $n$ 个根两两不同，对应的 Vandermonde 特征向量构成一组基，因此
\begin{equation}
 \det A=\prod_{\xi^n=\mu}f(\xi).
\end{equation}
该恒等式两端都是 $\mu,c_0,\ldots,c_{n-1}$ 的多项式，故也可由连续性推广到退化情形。
\end{solution}
